\section{ADH, fishing-mortality placement, and northern-cod records}\label{si:cross-system}

\subsection{ADH buffer-years recomputation}

The archived ADH measure is defined stock by stock as
\[
\mathrm{ADH}=\max\left\{0,\frac{\ln(\mathrm{SSB}/B_{\lim})}{F}\right\},
\]
where SSB is the current assessed spawning-stock biomass, \(B_{\lim}\) is the stock-specific limit reference point, and \(F\) is the current fishing mortality. It is an algebraic buffer-years proxy, not a direct measurement of institutional response or a guarantee of time to recovery.

The recovered input cohorts include an archived 43-stock table, a 454-stock RAM v4.66 broader pool, a 415-row v4.44 pool, and the screen-cohort summaries. Recalculation matches deposited ADH to numerical precision (maximum absolute difference \(1.14\times10^{-13}\) in the v4.66 broad pool; the archived 43-stock values match exactly to printed precision). Descriptive medians and zero fractions reproduce: 43-stock archived median 1.7902 years with 18.60\% zero values; v4.66 broad-pool median 3.3893 years with 15.20\% zero values; v4.66 forage-fish subset median 2.6962 years with 18.75\% zero values. In the 33-stock paired vintage comparison, the v4.44 and v4.66 medians are 2.7901 and 2.3442 years, respectively; across 396 paired broader stocks, median absolute ADH change is 0.7958 years. Data vintage and cohort definition therefore matter.

A fresh 20,000-draw seed-11 resampling run from the v4.66 broad pool yields median-of-pool draws with mean 3.3452 years and percentile range 1.8423--5.5307 years; 2.33\% of draws fall below the archived 43-stock median of 1.7902. The archived resampling record gives 2.12\%. The forage subset's new seed-11 run has median-of-pool draw mean 2.7916 years, percentile range 2.3081--3.5649, and 0.39\% below the archived median (archived value 0.29\%). This resampling describes a particular cohort-vintage construction; it is not a probability that a given stock is sustainable.

\subsection{Fishing-mortality data verification and placement}

The local F script first checks three gates. (G1) It reproduces the screen cohort: 42 stocks with at least 20 valid SSB observations. (G2) It compares fisheries-ADH F with both mean F and last F: the archived fisheries-ADH values match last F for 42 of 43 intersected stocks; \texttt{\detokenize{SPRBLKGSA29}} is the exception (archived fisheries-ADH F 0.2100, series mean F 0.2344, and last-F 0.0400, with 32 F observations). (G3) It reproduces the v4.66 broader-pool F quartiles (0.0800, 0.2143, 0.4639) and all 13 taxonomic-group quartile triples. These are arithmetic and source-cohort checks.

Among 38 of the 42 screen stocks with usable F series, the mean of stock-level F medians is 0.3663 (IQR 0.2701--0.7243; range 0.0053--2.6776). In the 454-stock broader pool, \(F_{\rm now}\) has median 0.2143 (IQR 0.0800--0.4639). The 1170 paired TC/SSB and F observations have Spearman \(\rho=0.8243\), \(p=7.88\times10^{-291}\), showing that the two exploitation summaries are strongly associated in this merged record. A comparison of model \(F^*=0.00209\) with empirical F values places most empirical values outside the illustrative model's low-F grid; because the model scale is not calibrated to these stocks, this is an explicit scale-mismatch warning rather than evidence that real fisheries are infeasible.

The exploratory rank association between stock-level mean F and target-band null-excess is positive but not conventionally significant in the screen cohort: band A \(\rho=0.3069, p=0.0609\), band B \(\rho=0.2017, p=0.2247\). Excluding the seven Peru/Chile cases reduces these to band A \(\rho=0.0339, p=0.8565\) and band B \(\rho=0.0327, p=0.8615\). The attenuation is consistent with region-linked confounding, but it does not identify the confound or establish a causal effect. The per-stock F series and model-placement labels are shown below.

\begin{landscape}
\scriptsize
\setlength{\tabcolsep}{5pt}
\begin{longtable}{@{}p{3.2cm}r r r r r l@{}}
\caption{Screen-cohort fishing-mortality series and placement categories. Values are retained as the local v53 output; a dash indicates no usable F series.}\label{tab:fplacement}\\
\toprule Stock ID & SSB \(n\) & F \(n\) & Mean F & Median F & Max F & Placement \\
\midrule\endfirsthead
\toprule Stock ID & SSB \(n\) & F \(n\) & Mean F & Median F & Max F & Placement \\
\midrule\endhead
\bottomrule\endfoot
\texttt{ANCHIXa} & 33 & 32 & 1.0000 & 0.9697 & 1.8305 & beyond-feasible \\
\texttt{ANCHMEDGSA17-18} & 48 & 48 & 0.6385 & 0.5540 & 1.5400 & beyond-feasible \\
\texttt{ANCHMEDGSA29} & 34 & 34 & 1.2613 & 1.1854 & 1.9805 & beyond-feasible \\
\texttt{ANCHMEDGSA7} & 20 & 20 & 0.1085 & 0.1060 & 0.2310 & beyond-feasible \\
\texttt{ANCHOBAYB} & 34 & 27 & 0.3336 & 0.3640 & 0.6810 & beyond-feasible \\
\texttt{ANCHOSA} & 36 & 0 & - & - & - & no-F \\
\texttt{ARGANCHONARG} & 30 & 30 & 0.0437 & 0.0420 & 0.0960 & beyond-feasible \\
\texttt{ARGANCHOSARG} & 30 & 30 & 0.0053 & 0.0040 & 0.0140 & unstable-range \\
\texttt{HERR2224IIIa} & 31 & 29 & 0.4220 & 0.4070 & 0.7010 & beyond-feasible \\
\texttt{HERR2529-32} & 48 & 47 & 0.2451 & 0.2190 & 0.4600 & beyond-feasible \\
\texttt{HERR2532} & 38 & 37 & 0.2701 & 0.2430 & 0.4820 & beyond-feasible \\
\texttt{HERR30} & 44 & 42 & 0.1411 & 0.1425 & 0.2060 & beyond-feasible \\
\texttt{HERR30-31} & 61 & 60 & 0.0841 & 0.0587 & 0.2564 & beyond-feasible \\
\texttt{HERR31} & 31 & 31 & 0.3294 & 0.2880 & 0.6400 & beyond-feasible \\
\texttt{HERR4RFA} & 53 & 30 & 0.0841 & 0.0705 & 0.1990 & beyond-feasible \\
\texttt{HERR4RSP} & 53 & 39 & 0.2128 & 0.1750 & 0.5780 & beyond-feasible \\
\texttt{HERR4TFA} & 41 & 36 & 0.3270 & 0.3275 & 0.9060 & beyond-feasible \\
\texttt{HERR4TSP} & 41 & 36 & 0.4930 & 0.4240 & 1.7400 & beyond-feasible \\
\texttt{HERR4VWX} & 42 & 41 & 0.7083 & 0.5700 & 2.4600 & beyond-feasible \\
\texttt{HERRIsum} & 37 & 31 & 0.2701 & 0.2600 & 0.4400 & beyond-feasible \\
\texttt{HERRNIRS} & 63 & 62 & 0.4352 & 0.3830 & 0.9930 & beyond-feasible \\
\texttt{HERRNORSS} & 74 & 73 & 0.2945 & 0.1500 & 3.4900 & beyond-feasible \\
\texttt{HERRNS-IIIa-VIId} & 77 & 76 & 0.3384 & 0.2520 & 1.2100 & beyond-feasible \\
\texttt{HERRNWATLC} & 58 & 57 & 0.3777 & 0.2750 & 0.7980 & beyond-feasible \\
\texttt{HERRRIGA} & 47 & 45 & 0.4185 & 0.4240 & 0.7080 & beyond-feasible \\
\texttt{HERRSIRS} & 64 & 63 & 0.4734 & 0.4500 & 0.9740 & beyond-feasible \\
\texttt{HERRVIa} & 57 & 57 & 0.3549 & 0.2720 & 1.0800 & beyond-feasible \\
\texttt{HERRVIaVIIbc} & 62 & 61 & 0.3409 & 0.2760 & 1.0400 & beyond-feasible \\
\texttt{JANCHOPJPN} & 36 & 0 & - & - & - & no-F \\
\texttt{JANCHOSETO} & 38 & 0 & - & - & - & no-F \\
\texttt{JANCHOTSST} & 42 & 0 & - & - & - & no-F \\
\texttt{PANCHCCH} & 36 & 36 & 0.5616 & 0.5025 & 1.7000 & beyond-feasible \\
\texttt{PANCHCHVX} & 23 & 23 & 1.3357 & 1.3700 & 3.1100 & beyond-feasible \\
\texttt{PANCHCSCH} & 31 & 31 & 1.2472 & 1.0500 & 3.6200 & beyond-feasible \\
\texttt{PANCHNCH} & 29 & 15 & 2.6776 & 2.9800 & 4.0200 & beyond-feasible \\
\texttt{PANCHNCHSP} & 36 & 36 & 2.6179 & 1.9250 & 7.9800 & beyond-feasible \\
\texttt{PANCHPERUNC} & 47 & 47 & 0.7297 & 0.6150 & 2.1000 & beyond-feasible \\
\texttt{PANCHSCH} & 28 & 28 & 0.8514 & 0.7800 & 2.7800 & beyond-feasible \\
\texttt{SPRAT22-32} & 50 & 48 & 0.3179 & 0.3189 & 0.5320 & beyond-feasible \\
\texttt{SPRATIIIa-IV} & 48 & 47 & 1.2363 & 1.2750 & 2.2030 & beyond-feasible \\
\texttt{SPRATNS} & 45 & 44 & 1.0598 & 1.0185 & 2.3200 & beyond-feasible \\
\texttt{SPRBLKGSA29} & 30 & 32 & 0.2344 & 0.2050 & 0.6400 & beyond-feasible \\
\end{longtable}
\end{landscape}


\subsection{Northern-cod assessment table}

\begin{table}[htbp]
\centering\small
\caption{Northern-cod assessment values shown in the descriptive two-window figure. The reported instantaneous natural mortality \(M\) comes from the assessment table; \(e^{-M}\) is a transformation, not a separate observed survival series.}
\label{tab:cod-values}
\begin{tabular}{@{}rrrr@{}}
\toprule Year & SSB (kt) & \(M\;(\mathrm{yr}^{-1})\) & \(e^{-M}\) \\
\midrule
1991 & 735.00 & 1.002 & 0.367 \\
1992 & 382.00 & 2.214 & 0.109 \\
1993 & 101.00 & 2.575 & 0.076 \\
1994 & 31.00 & 2.331 & 0.097 \\
1995 & 10.00 & 0.288 & 0.750 \\
1996 & 16.05 & 0.341 & 0.711 \\
2000 & 34.42 & 0.717 & 0.488 \\
2004 & 20.07 & 0.362 & 0.696 \\
2005 & 25.18 & 0.288 & 0.750 \\
2010 & 96.91 & 0.696 & 0.499 \\
2015 & 298.65 & 0.278 & 0.757 \\
\bottomrule
\end{tabular}
\end{table}

The values were cross-checked against the DFO CSAS SAR 2016/026 table. The crash window is 1991--1995; the later plotted interval extends to 2015 and includes low-stock persistence and subsequent increases. The plot partitions the chronology but does not decompose causation. No numerical constrained-\(M\) allocation is included in v53: the raw source quantities and a fully reproducible calculation were not recovered. Claims about the size of a seal-attributable, fishing-attributable, or governance-attributable component are therefore not made.

